\section{Related Work and Limitations}
\label{sec:related}

Prior research optimizes fixed-budget grouped error, expensive predicate execution, online stopping, model ranking, or single-population prediction-assisted inference.
\textsc{StratQ} brings these threads into a unified fixed-snapshot SQL grouping contract where every group shares a single anytime-valid coverage event that directly certifies \texttt{COUNT}, \texttt{HAVING}, and \texttt{TOP-$k$} outputs or abstains.

\subsection{Grouped Approximate Query Processing}
Grouped AQP has a rich database lineage. Congressional Sampling~\cite{acharya2000congress} allocates offline stratified samples via Senate/House designs; Sample+Seek~\cite{sample2016}, Quickr~\cite{kandula2016quickr}, and ShadowAQP~\cite{shadow2023} add distribution precision, inline/universe sampling, and attribute-oriented allocation/generation. BlinkDB~\cite{agarwal2013blinkdb}, AQP++~\cite{peng2018aqpplus}, VerdictDB~\cite{park2018verdict}, LWS~\cite{walenz2019lws}, ApproxHadoop~\cite{goiri2015approxhadoop}, and ApproxIoT~\cite{wen2018approxiot} cover multi-resolution samples, aggregate cubes, variational subsampling, learned sampling, and MapReduce/edge analytics.
PilotDB~\cite{zhu2025pilotdb} formalizes a-priori simultaneous relative error across groups. These systems target fixed-budget or two-stage CLT bounds. PilotDB excludes aggregates used as \texttt{GROUP BY} keys (e.g., \texttt{GROUP BY COUNT(*)}), not ordinary \texttt{COUNT} outputs. \textsc{StratQ} instead governs every look, route, and budget cap with one finite-sample anytime event.

\subsection{Expensive Predicates and Semantic Engines}
Expensive predicates are an active VLDB/SIGMOD frontier. NoScope~\cite{kang2017noscope}, Probabilistic Predicates~\cite{lu2018pp}, CORE~\cite{yang2022core}, and TASTI~\cite{kang2022tasti} reduce inference through filters, proxies, and semantic indexes without certifying population totals. SUPG~\cite{kang2020supg} provides asymptotic selection precision/recall; BlazeIt~\cite{kang2019blazeit} uses neural control variates with always-valid empirical-Bernstein stopping; InQuest~\cite{russo2023inquest} gives expected $O(1/B)$ error under stationary streams.
LOTUS~\cite{patel2025lotus} provides gold-relative semantic accuracy with fixed-sample CLT/multiplicity corrections; UQE~\cite{dai2024uqe} discovers semantic groups; Palimpzest~\cite{liu2025palimpzest} and DocETL~\cite{shankar2025docetl} optimize AI pipelines; ThalamusDB~\cite{jo2024thalamusdb} emits anytime partial results with deterministic bounds. BARGAIN~\cite{zeighami2026bargain} targets a single global cascade quality metric; Conformal Cascade~\cite{dou2026conformalcascade} gives marginal acceptance bounds ($1-K\alpha$, or $1-\alpha$ under selection preservation). \textsc{StratQ} is orthogonal to proxy generation and adds a per-group SQL count contract.

\subsection{Online Aggregation, Stopping, and Ranking}
Online Aggregation~\cite{hellerstein1997online}, Ripple Joins~\cite{haas1999ripple}, Wander Join~\cite{li2016wander}, and DeepOLA~\cite{sheoran2023deepola} provide progressive estimates and running intervals for scans/joins, not fixed-population cross-group certification. Anytime top-$k$ measures~\cite{arai2009anytimetopk} and bandits---AT-LUCB~\cite{jun2016atlucb}, Thresholding Bandits~\cite{tao2019threshold}, and APGAI~\cite{jourdan2026apgai}---target i.i.d.-arm ranking error, regret, or good-arm identification rather than fixed group counts.
Concurrent Rank CS~\cite{khosravi2026rankcs}, SERPANT~\cite{gu2026serpant}, and Sequential Model Confidence Sets~\cite{arnold2026smcs} already provide anytime leaderboards, FWER control, and best-model subsets. Thus a joint anytime top-$k$ event is not itself novel; \textsc{StratQ}'s distinction is fixed-snapshot SQL group counts, within-group WoR, and relational execution cost.

\subsection{Finite-Population Sequential Inference}
\textsc{StratQ} builds on SAVI~\cite{ramdas2023savi}, stitched/empirical-Bernstein CS~\cite{howard2021cs}, WoR betting~\cite{wor2020,waudby2024betting}, universal-portfolio CS~\cite{orabona2024portfolio}, and risk-limiting audits~\cite{waudby2021rilacs}; tighter fixed-$n$ WoR rate bounds~\cite{shekhar2026rates} are non-sequential. Turner and Gr\"unwald~\cite{turner2023stratified} study cross-stratum e-variables.
The closest preprints are Spertus et al.~\cite{spertus2026sequential}, with predictable cross-stratum selection, WoR remaining means, and simultaneous LCBs (union-of-intersections sequences need not share one global test supermartingale), and Hait~\cite{hait2026target}, with adaptive allocation, target-set coverage, and best-system identification under i.i.d.\ within-stratum observations. \textsc{StratQ} adopts the WSR engine; its statistical contribution is predictable batch-refitted row-level residual control variates plus Safe-LOO gating in one global WoR filtration, with directed binary64 rounding.

\subsection{Prediction-Powered and Active Inference}
Prediction-Powered Inference~\cite{angelopoulos2023ppi} (including finite-population WoR betting in its appendix), PPI++~\cite{angelopoulos2026ppipp}, Cross-PPI~\cite{zrnic2024crossppi}, and Stratified PPI~\cite{fisch2024stratppi} accelerate fixed-sample estimation. Active Statistical Inference~\cite{zrnic2024active} permits predictable refitting; Filtered PPI~\cite{xu2026fppi} motivates selective use of noisy predictions but its main efficiency result is asymptotic.
Prediction-Powered E-Values~\cite{csillag2025evalues} give exact time-uniform PPI, Kilian et al.~\cite{kilian2025anytime} give asymptotic anytime PPI, Sfyraki and Wang~\cite{sfyraki2026active} study active sequential mean estimation, and FAQ~\cite{wu2026faq} combines online refitting with proactive querying (its main martingale CLT requires with-replacement sampling). CELEUS~\cite{zhou2026celeus} uses fixed surrogate completion and IPW residuals under WoR for a single risk metric. Proposition~\ref{prop:celeus-equiv} shows our fixed-proxy uniform-WoR specialization shares its unbiasedness identity; we extend to predictable Safe-LOO refitting, simultaneous grouped SQL contracts, and empirical operator comparisons.

\subsection{Control Variates and Analytical Execution}
Model-assisted survey sampling---difference/GREG estimators~\cite{sarndal1992model} and ML-assisted extensions~\cite{breidt2017modelassisted}---provides the design-based control-variate foundation, typically under single-look inference. DuckDB~\cite{raasveldt2019duckdb} and Arrow~\cite{apachearrow} supply vectorized projection and columnar transfer. Our evaluation therefore reports calls, latency, memory, and certified decisions separately rather than treating proxy accuracy as query validity.

\subsection{Limitations}
\label{sec:limitations}

\textsc{StratQ} carries clear scope boundaries that guide deployment and future work:
(1)~\emph{Static snapshot and known group keys}: the current contract assumes a fixed table with cheap, known group keys; streaming updates, concept drift, and unknown group keys produced by an expensive model require separate mechanisms.
(2)~\emph{Per-query risk budgets}: each query consumes its own $\delta$; interactive sessions issuing many adaptive follow-up queries over the same snapshot require online FDR/FWER risk accounting across queries.
(3)~\emph{Workload-dependent proxy gains}: as shown on Adult ($+0.93\%$) and Yelp \texttt{COUNT}, weak proxies or strongly pre-stratified groups yield little or zero variance reduction beyond uniform WoR. Safe-LOO is a predictable cost gate, not a worst-case call guarantee; the Yelp gate trace shows a 12.9\% mean LOO variance reduction but a $12.7\%$ realized call overhead vs.\ Spertus-WoR, with negligible prediction range.
(4)~\emph{Profile materialization, batching, and memory attribution}: at 10M rows, offline-replay columnar profile construction dominates peak memory (0.68\,GB end RSS) and needs a modeled $Q^*=2$ queries to amortize against cold runs. The query-stage $+21.69$\,MB RSS delta is one immediate pre-teardown observation; its traced-allocation peak is 85.6\,MB and final post-teardown delta is $-0.68$\,MB. Sub-second CPU predicates (such as SMS $k$-NN) still favor exhaustive batching.
(5)~\emph{Floating-point assumptions}: Theorem~\ref{thm:fp} assumes finite binary64 inputs and valid outward feature bounds supplied by the safe engine path; formal machine-checked proof verification (e.g., Lean/Coq) remains future work.
(6)~\emph{Finite empirical trials}: observed coverage across 195 paired Sprint~1 runs, 30 Zipf runs over 15 configurations, and 960 decision-boundary trials corroborates the theorems; it does not replace proof or establish population coverage. Three-seed GPU and proxy-sweep results are descriptive, and two-seed Zipf cells are especially limited.

\section{Conclusion}
\label{sec:conclusion}

\textsc{StratQ} reframes proxy-assisted analytical querying around an explicit, anytime-valid grouped SQL execution contract.
By combining uniform within-group WoR sampling, predictable Safe-LOO residual control variates, cross-group adaptive routing, and directed binary64 outward rounding under a single $(1-\delta)$ simultaneous coverage event, \textsc{StratQ} unifies grouped \texttt{COUNT} estimation, \texttt{HAVING} classification, and \texttt{TOP-$k$} membership certification with explicit abstention.
Implemented over DuckDB and Apache Arrow and evaluated on live GPU neural predicates, 10M-row Parquet snapshots, eleven database/modern operator adaptations plus an independent hypergeometric CS, \textsc{StratQ} delivers $3.61\times$--$4.78\times$ ratio-of-medians query-level GPU speedups after proxy scoring, demonstrates the value of decision-aware routing over post-hoc \texttt{COUNT}-derived decisions on tight boundaries, and reports profile-dominated 10M-row memory with an 85.6\,MB query traced-allocation peak rather than a retained-RSS guarantee.
